\documentclass[../../main.tex]{subfiles}
\begin{document}

{\bf [解]}

$f(n)$は $n!$の中の素因数$2^p,5^q\ (p,q\in\mathbb{Z}_{\ge0})$について，$p,q$のうち大きくない方に等しい．明らかに$p\ge q$だから$f(n)=q$である．

\paragraph{(1)}

まず$n=10$について$10\div5=2$だから，$q=2\ \therefore\ f(10)=2$である．
次に$n=100$について$\dfrac{100}{5}=20$，$\dfrac{100}{25}=4$だから，$f(100)=20+4=24$である．

\paragraph{(2)}

(1)と同様にして$f(10^n)$は
\begin{align}
f(10^n)=\sum_{k=1}^\infty\left[\dfrac{10^n}{5^k}\right] \quad\cdots\text{①} \label{eq:1}
\end{align}
とかける．ただしガウス記号$[x]$は$x$以下の最大の整数である．$k$についての和は無限大までとるが，実際には$k=2n$の段階で
$k=2n$の時，
\begin{align}
\left[\dfrac{10^n}{5^{2n}}\right]=\left[\left(\dfrac{2}{5}\right)^n\right]=0
\end{align}
となるから，$[10^n/5^k]$が非負で$k$の単調減少関数であることとあわせて，$k \ge 2n$では全ての寄与は$0$である．従って和は$k=1,2,\cdots 2n$までに制限してよく，
\begin{align}
  f(10^n)=\sum_{k=1}^{2n}\left[\dfrac{10^n}{5^k}\right] \quad\cdots\text{②} \label{eq:2}
\end{align}
である．

ここで$[x]$の性質から
\begin{align}
\dfrac{10^n}{5^k}-1<\left[\dfrac{10^n}{5^k}\right]\le\dfrac{10^n}{5^k} \label{eq:3}
\end{align}
だから，$k=1,2,\cdots,2n$として足し合わせて\cref{eq:2}を代入すると
\begin{align}
\begin{split}
\sum_{k=1}^{2n}\left(\dfrac{10^n}{5^k}-1\right)&<f(10^n)\le\sum_{k=1}^{2n}\dfrac{10^n}{5^k} \\
10^n\cdot\dfrac{1}{5}\cdot\dfrac{1-(1/5)^{2n}}{1-1/5}-2n&<f(10^n)\le10^n\cdot\dfrac{1}{5}\cdot\dfrac{1-(1/5)^{2n}}{1-1/5} \\
10^n\dfrac{1}{4}\left(1-(1/5)^{2n}\right)-2n&<f(10^n)\le 10^n\dfrac{1}{4}\left(1-(1/5)^{2n}\right)
\end{split}
\end{align}
である．

両辺$10^n$でわって
\begin{align}
 \dfrac{1}{4}\left(1-(1/5)^{2n}\right) -\dfrac{2n}{10^n}<\dfrac{f(10^n)}{10^n}\le \dfrac{1}{4}\left(1-(1/5)^{2n}\right)
\end{align}
この両辺は共に$n\to\infty$で$\dfrac{1}{4}$に収束するから，はさみうちの定理から
\begin{align}
\dfrac{f(10^n)}{10^n}\longrightarrow\dfrac{1}{4}
\end{align}
である．

\newpage

\end{document}
